% Section 3: Forecasting for reinforcement learning
\subsection{Forecasting for reinforcement learning}
\label{sec:forecasting_for_rl}

\subsubsection{World models}
\label{sec:world_models}

The dominant paradigm for using forecasting inside RL is the world model. The
architecture was introduced in its modern form by \citet{HaSchmidhuber2018}. A
vision module $V$ compresses high-dimensional observations into a low-dimensional
latent code $z$, a memory module $M$ implemented as a Mixture Density Network
recurrent network predicts the distribution of $z_{t+1}$ given $z_t$ and $a_t$, and
a controller $C$ implemented as a simple linear policy chooses actions in the
latent space. The controller is trained by evolution strategies, and the entire
agent can be trained partially or wholly inside a learned dream of the world. On
VizDoom, an agent trained only inside the dream reaches roughly $80\%$ of the
performance of an agent trained on the actual environment.

The Dreamer family of \citet{Hafner2025} refines the architecture into the Recurrent
State Space Model, which combines a deterministic hidden state $h_t$ with a
stochastic latent $z_t$ and trains the world model jointly with an actor and a
critic that operate on imagined rollouts. The DreamerV3 release reports that a
single hyperparameter configuration solves more than $150$ benchmark tasks across
visual, proprioceptive, and discrete-state domains, and that the agent collects
diamonds in Minecraft from scratch without a curriculum. The advance over the prior
Dreamer versions is primarily engineering, namely layer normalisation, balanced
loss weighting that prevents the reconstruction term from dominating, and
activation normalisation that prevents value collapse. The methodological message
for our chapter is that learned probabilistic forecasters of dynamics can replace
hand-built simulators across a wide range of domains, including the ones in
which classical model-free RL would struggle.

\subsubsection{Trajectory generative models}
\label{sec:trajectory_models}

A parallel line treats the trajectory itself as a sequence to be modelled. The
Decision Transformer of \citet{Chen2021DT} recasts offline RL as conditional sequence
modelling. A causally masked transformer is trained on tuples of state, action, and
return-to-go, and at inference one conditions on a desired return $R^\star$ and
samples actions autoregressively. The Trajectory Transformer of \citet{Janner2021TT}
augments this with explicit beam search over discretised state-action-reward
sequences, and Diffuser \citep{Janner2022diffuser} replaces autoregressive
generation with iterative denoising of an entire trajectory under classifier
guidance for goal-conditioning. Decision Diffuser \citep{Ajay2023} pushes the idea
further, parametrising the policy as a return-conditional diffusion model and
dispensing with dynamic programming altogether. Across these models, the common
move is to treat planning as conditional generation of high-dimensional structured
objects, namely entire trajectories, with the forecast distribution serving as the
substrate for both policy evaluation and policy improvement.

\subsubsection{Counterfactual forecasting for off-policy evaluation}
\label{sec:counterfactual_ope}

Off-policy evaluation in finite samples is bias-variance constrained. Importance
sampling estimators are unbiased but high variance, and model-based estimators are
low variance but biased by model misspecification. A third route, due to
\citet{Buesing2019}, uses a structural causal model on the trajectory and produces
counterfactual estimates of the outcome that would have obtained under an
alternative policy. Given a logged transition $(s_t, a_t^{\mathrm{obs}},
s_{t+1}^{\mathrm{obs}})$ generated under an SCM
$s_{t+1} = f(s_t, a_t, u_t)$ with exogenous noise $u_t$, the noise $u_t$ is inferred
from the observed transition, and the alternative action $a_t^{\mathrm{cf}}$ is then
applied with the same noise realisation to yield the counterfactual next state.
The advantage relative to model-based OPE is that the residual noise is shared,
which removes a source of model error.

\citet{Oberst2019} formalise the discrete-action case via Gumbel-max structural
equations, and apply the framework to a sepsis management POMDP. \citet{TangWiens2023}
combine SCM-based counterfactual outcomes with importance sampling to obtain
variance reductions while preserving unbiasedness under correct human annotation. The
shared identifying assumption across this strand is that the SCM is correctly
specified, which is the precise place where the Lucas critique applies and is the
hardest open issue in the broader programme.

\subsubsection{Time series foundation models}
\label{sec:tsfm}

The most recent development is the arrival of pretrained univariate time series
foundation models. Chronos of \citet{Ansari2024} tokenises a real-valued time series
by mean-scaling and uniform quantisation into a fixed vocabulary of $4{,}096$ bins,
and trains a T5 family transformer ranging from $20$M to $710$M parameters on a
cross-entropy objective. Inference produces a categorical distribution over the
next bin, which is then mapped back to a real-valued quantile forecast by inverse
quantisation. TimesFM of \citet{Das2024} adopts a decoder-only architecture of $200$M
parameters and operates on patches of consecutive time points, pretrained on a
corpus of roughly $100$ billion time-point observations. Lag-Llama of
\citet{Rasul2023} eschews tokenisation in favour of explicit lag covariates, which
maps directly onto the classical autoregressive structure familiar to time series
econometrics.

All three models report zero-shot accuracy comparable to specialised models on
standard public benchmarks, and the natural question for our purposes is whether
they can serve as world models for sequential decision problems. The answer is
not yet, for three reasons. First, \citet{PerezDiaz2025} show that the standard
TSFM training objective optimises marginal calibration but not joint calibration
across horizons. The papers in their survey report low CRPS on each individual
horizon and high Energy Score on joint trajectory samples, and roughly two thirds
of recent TSFMs produce only point or parametric outputs that cannot answer
path-dependent operational queries. Second, \citet{Rahimikia2025} document that
on daily excess financial returns, zero-shot Chronos and TimesFM attain Sharpe
ratios near $0.05$, fine-tuning on financial data raises the Sharpe to roughly
$0.4$, and pretraining from scratch on financial data raises it further to between
$0.5$ and $0.6$. Third, \citet{Jetwiriyanon2025} examine macroeconomic indicators
across structural breaks (2008, 2020) and find that TSFM accuracy degrades during
the break but recovers faster than that of ARIMA and GARCH, with recovery times
of five to ten quarters. The collective picture is that TSFMs interpolate well in
distribution and degrade gracefully out of distribution, which is consistent with
their use as priors but inconsistent with their use as the action-conditional
dynamics model an RL agent would need.

The frontier on this point is \citet{Ada2025}. They train an offline RL policy on
non-stationary observation processes, using Lag-Llama in zero-shot mode to produce
a per-episode forecast of additive observation offsets and a diffusion model to
refine the within-episode belief. Their setup is the closest thing in the
literature to a TSFM-in-the-loop RL system, but the TSFM is a passive
observation-process prior rather than an action-conditional dynamics model. The
authors note that zero-shot foundation-model accuracy degrades when real-world
offsets violate the predefined assumptions, and that partial identifiability,
namely the inability to disentangle policy effect from exogenous drift from a
single observation, motivates their diffusion-fusion step. The action-conditional
foundation forecaster remains the next architectural step.
